\section{Real photographs}
\label{sec:real}

The reviews note that the real-image experiment uses one dataset, reports medians without uncertainty, restricts the pairs, and includes no geometric baseline. This section gives the complete protocol, including one input that the main paper does not mention, and the results on all sampled pairs.

\noindent\textbf{Data.} The \rlImages{} photographs used come from flights of a DJI Phantom~4 RTK over a village, with the camera tilted by about \rlTilt$^\circ$ from the nadir and poses from RTK positioning and bundle adjustment. The photographs are undistorted and resampled to the calibration $K$ of PairUAV, and the labels are computed from the poses as for PairUAV. Among all ordered pairs whose cameras are at most 30\,m apart we sample 600: \rlNin{} with a baseline of at most \rlInMax\,m, which is the baseline-to-height range covered by PairUAV once scaled to this flight height (38\,m at 220\,m), and \rlNout{} beyond it. Table~4 of the main paper reports the first group, as its caption states; the second group was sampled at the same time and is reported here.

\noindent\textbf{Flight height.} DreamNav is not fine-tuned and sees no label of this dataset. It does, however, receive one number: the height above ground reported by the aircraft, 60\,m. The two DreamNav rows of Table~4 multiply the predicted displacement by $60/220.2$, the ratio of this height to the flight height at which the scale cues were calibrated (\cref{sec:cands}). The main paper omits this, and the baselines of Table~4 do not receive it. Without the flight height the output is too long by this factor and the median endpoint error is \rlNoHeightPos\,m (\cref{tab:real}); giving the same factor to the regression on MASt3R statistics lowers its median error to \rlMlpPos\,m. A height above ground is available on most UAVs (barometer, rangefinder or terrain model), but the claim of Sec.~5.3 should read ``zero-shot given the flight height''.

\begin{table*}[t]
\centering
\caption{\textbf{Real photographs: Table~4 re-scored with uncertainty}, in the camera frame of the first image unless noted, for the pairs within the baseline-to-height range of the training data (the pairs of Table~4) and beyond it. EPE: endpoint error; Table~4 reports its median as ``Pos.'' and \SRd{} as SR. Brackets: 95\% bootstrap intervals over source images. Angular errors of $\dth$ are mean / median. ``No flight height'': the output is not rescaled by the height reported by the aircraft.}
\label{tab:real}
\footnotesize
\setlength{\tabcolsep}{4pt}
\resizebox{\textwidth}{!}{%
\begin{tabular}{@{}lccccccccc@{}}
\toprule
& \multicolumn{2}{c}{EPE (m)} & \SRe@1\,m & \SRe@2\,m & \SRd@1\,m & \SRd@2\,m & $\dth$ error ($^\circ$) & $\dpsi$ MAE & Dir. med. \\
Method & median & mean & (\%) & (\%) & (\%) & (\%) & mean / med. & ($^\circ$) & ($^\circ$) \\
\midrule
\multicolumn{10}{@{}l}{\emph{Within the training range: 300 pairs, baselines 0.7--10.3\,m}}\\
ViNT & 7.75 \ci{7.55}{8.05} & 7.32 & 0.3 \ci{0.0}{1.0} & 3.3 & 3.7 & 8.0 & 71.34 / 60.29 & 14.53 & 90.6 \\
NoMaD & 21.21 \ci{20.31}{21.95} & 20.76 & 0.0 \ci{0.0}{0.0} & 0.0 & 0.3 & 2.3 & 86.01 / 81.75 & 17.86 & 89.5 \\
$\pi_0$-FAST & 25.93 \ci{25.46}{26.51} & 26.06 & 0.0 \ci{0.0}{0.0} & 0.0 & 0.0 & 0.0 & 94.26 / 113.22 & 4.43 & 91.8 \\
Constant (zero motion) & 7.84 \ci{7.45}{8.08} & 7.18 & 1.3 \ci{0.3}{2.7} & 3.7 & 1.3 & 3.7 & 83.71 / 89.34 & 0.45 & 90.0 \\
MASt3R-statistics MLP, flight height & 5.10 \ci{4.79}{5.30} & 5.28 & 1.7 \ci{0.3}{3.2} & 8.3 & 18.3 & 34.3 & 3.72 / 1.93 & 3.40 & 41.0 \\
\cmidrule(l){1-10}
Ours (Stage I), no flight height & 15.18 \ci{14.23}{16.09} & 14.95 & 0.0 \ci{0.0}{0.0} & 0.0 & 2.3 & 3.3 & 2.31 / 0.72 & 1.90 & 4.4 \\
Ours (Stage II), no flight height & 16.42 \ci{15.64}{17.58} & 16.39 & 0.0 \ci{0.0}{0.0} & 0.0 & 2.0 & 3.7 & 2.25 / 0.67 & 1.62 & 1.9 \\
Ours (Stage I) & 1.43 \ci{1.25}{1.77} & 2.74 & 34.0 \ci{29.0}{39.4} & 59.0 & 45.0 & 69.3 & 2.31 / 0.72 & 1.90 & 4.4 \\
Ours (Stage II) & 0.84 \ci{0.72}{0.98} & 2.13 & 57.7 \ci{51.5}{63.8} & 72.3 & 63.3 & 80.3 & 2.25 / 0.67 & 1.62 & 1.9 \\
Ours (Stage II), ENU with learned heading & 4.85 \ci{4.24}{5.57} & 5.43 & 7.3 \ci{4.5}{10.5} & 18.0 & 63.3 & 80.3 & 2.25 / 0.67 & 1.62 & 43.0 \\
\midrule
\multicolumn{10}{@{}l}{\emph{Beyond the training range: 300 pairs, baselines 10.4--29.8\,m}}\\
Constant (zero motion) & 21.71 \ci{20.50}{22.80} & 20.81 & 0.0 \ci{0.0}{0.0} & 0.0 & 0.0 & 0.0 & 101.73 / 91.45 & 0.55 & 90.0 \\
MASt3R-statistics MLP, flight height & 18.88 \ci{18.03}{19.48} & 17.81 & 0.0 \ci{0.0}{0.0} & 0.0 & 0.0 & 2.0 & 19.95 / 6.52 & 7.27 & 55.5 \\
\cmidrule(l){1-10}
Ours (Stage I), no flight height & 22.29 \ci{21.57}{22.87} & 21.84 & 0.0 \ci{0.0}{0.0} & 0.0 & 1.0 & 2.7 & 17.37 / 3.88 & 5.12 & 48.1 \\
Ours (Stage II), no flight height & 24.13 \ci{23.34}{24.83} & 24.59 & 0.0 \ci{0.0}{0.0} & 0.0 & 1.0 & 2.3 & 17.11 / 3.32 & 4.94 & 45.5 \\
Ours (Stage I) & 18.27 \ci{17.13}{19.33} & 16.59 & 0.0 \ci{0.0}{0.0} & 4.3 & 0.0 & 7.3 & 17.37 / 3.88 & 5.12 & 48.1 \\
Ours (Stage II) & 18.03 \ci{16.45}{19.01} & 15.92 & 1.7 \ci{0.3}{3.2} & 9.0 & 5.0 & 11.0 & 17.11 / 3.32 & 4.94 & 45.5 \\
Ours (Stage II), ENU with learned heading & 20.06 \ci{19.14}{20.94} & 18.82 & 0.0 \ci{0.0}{0.0} & 1.0 & 5.0 & 11.0 & 17.11 / 3.32 & 4.94 & 65.0 \\
\bottomrule
\end{tabular}}
\end{table*}


\noindent\textbf{Within the training range.} After refinement the median endpoint error is \rlSiiPos\,m and the endpoint success \rlSiiSRe\% \ci{\rlSiiSReLo}{\rlSiiSReHi}, against \rlSiPos\,m and \rlSiSRe\% \ci{\rlSiSReLo}{\rlSiSReHi} for the coarse stage; the paired gain is +\rlGainSRe{} points \ci{\rlGainSReLo}{\rlGainSReHi}. The mean endpoint error (\rlSiiEPEmean\,m) is well above the median: a minority of pairs fails by a large margin, which a median hides. The gain of the refinement on photographs with 3.7 times the relative relief of PairUAV supports the analysis of \cref{sec:planarity}.

\noindent\textbf{Beyond the training range.} For baselines between \rlInMax{} and \rlOutMax\,m DreamNav fails: the median endpoint error is \rlSiiOutPos\,m, close to the \rlZeroOutPos\,m of predicting no motion, and the endpoint success is \rlSiiOutSRe\%. The rotation is still estimated to a median of \rlSiiOutThetaMed$^\circ$ in heading, but the translation direction is lost (median error \rlSiiOutDir$^\circ$). DreamNav does not extrapolate beyond the baseline-to-height ratio of its training data, and the statement of Sec.~5.3 that it ``generalizes well'' holds only inside it.

\noindent\textbf{Absolute heading.} All numbers above are in the camera frame. Rotating the output into ENU with the learned heading of \cref{sec:absh} reduces the endpoint success within range from \rlSiiSRe\% to \rlEnuSRe\%: on these photographs the heading prior is wrong by \rlEnuHead$^\circ$ in the median (\cref{sec:heading}). On a real aircraft the heading must come from the compass.

\noindent\textbf{Geometric baselines.} \Cref{tab:real_geo} applies the pipelines of \cref{sec:geometric} to the same pairs. Nothing is fitted on the photographs: the inlier thresholds are those selected on PairUAV and the flight-height prior uses $H{=}60$\,m, the number that our own rows receive. Within the training range the SIFT pipeline has a median endpoint error of \grSiftEPEmed\,m \ci{\grSiftEPELo}{\grSiftEPEHi} and an endpoint success of \grSiftSRe\% \ci{\grSiftSReLo}{\grSiftSReHi}, which is \grSiftVsII{} points \ci{\grSiftVsIILo}{\grSiftVsIIHi} \emph{above} our refined stage (\grSiiSRe\%). RoMa, on the \grRomaN{} pairs on which it was run, has a median error of \grRomaEPEmed\,m and an endpoint success of \grRomaSRe\% (\grRomaVsII{} points \ci{\grRomaVsIILo}{\grRomaVsIIHi} above the refined stage on these pairs). LoFTR (\grLoftrSRe\%) and DeDoDe (\grDedodeSRe\%) with the rotation search are level with the refined stage, and our coarse stage (\grSiSRe\%) is below every keypoint pipeline of the table. Beyond the training range the SIFT pipeline finds no pose for \groSiftFail\% of the pairs and succeeds on \groSiftSRe\%, against \groSiiSRe\% for the refined stage, which returns an estimate for every pair. On real photographs DreamNav therefore does not improve over classical two-view geometry given the same flight height. We consider this the most important limitation that the additional experiments reveal: the learned fusion, which provides the advantage in distance on PairUAV (\cref{sec:georesults}), was fitted on rendered views at one flight height and does not carry that advantage to this dataset.

\begin{table*}[t]
\centering
\caption{\textbf{Geometric baselines on the real photographs}, zero-shot, with the flight height of 60\,m given to every method (our rows: displacement rescaled by $60/220$; baselines: flight-height prior with $H{=}60$\,m). Camera frame; 95\,\% intervals by bootstrap over source photographs. Last column: difference of \SRe{} to our refined stage on the same pairs.}
\label{tab:real_geo}
\footnotesize
\setlength{\tabcolsep}{4pt}
\resizebox{\textwidth}{!}{%
\begin{tabular}{@{}llrccccccc@{}}
\toprule
& & & Fail & \SRd@1\,m & EPE med. & \SRe@1\,m & $\dth$ med. & Dir. med. & $\Delta$\SRe{} \vs \\
Method & Rot. & Pairs & (\%) & (\%) & (m) & (\%) & ($^\circ$) & ($^\circ$) & Stage~II (points) \\
\midrule
\multicolumn{10}{@{}l}{\emph{Within the training range} (mean true distance 7.2\,m)}\\
SIFT + ratio test & inv. & 300 & 3.3 & 76.0 \ci{71.0}{80.5} & 0.53 \ci{0.43}{0.60} & 72.7 \ci{67.4}{77.4} & 0.549 & 1.9 & $+15.0$ \ci{+8.2}{+21.1} \\
SuperPoint + SuperGlue, our first pass & ours & 300 & 4.3 & 52.3 \ci{45.9}{58.2} & 1.32 \ci{0.96}{1.82} & 44.7 \ci{38.7}{50.3} & 0.845 & 4.9 & $-13.0$ \ci{-19.0}{-7.3} \\
SuperPoint + LightGlue & TTA & 300 & 8.7 & 46.0 \ci{40.5}{51.5} & 1.67 \ci{1.33}{2.16} & 36.7 \ci{31.4}{42.1} & 0.886 & 8.0 & $-21.0$ \ci{-27.5}{-14.8} \\
ALIKED + LightGlue & TTA & 300 & 5.0 & 59.0 \ci{53.7}{64.5} & 1.06 \ci{0.89}{1.40} & 48.3 \ci{42.7}{54.1} & 0.749 & 4.7 & $-9.3$ \ci{-15.4}{-3.3} \\
LoFTR & TTA & 300 & 1.0 & 66.3 \ci{61.2}{71.8} & 0.64 \ci{0.49}{0.74} & 61.0 \ci{55.4}{66.4} & 0.682 & 2.5 & $+3.3$ \ci{-2.3}{+9.0} \\
DeDoDe + rotation steerer & eq.+TTA & 300 & 2.3 & 65.0 \ci{58.9}{70.8} & 0.75 \ci{0.62}{0.93} & 58.0 \ci{51.9}{64.0} & 0.818 & 3.0 & $+0.3$ \ci{-5.3}{+5.9} \\
RoMa & TTA & 80 & 3.7 & 76.2 \ci{66.7}{85.7} & 0.33 \ci{0.27}{0.50} & 75.0 \ci{64.9}{84.7} & 0.784 & 1.3 & $+21.3$ \ci{+10.8}{+31.3} \\
MASt3R descriptors & TTA & 100 & 2.0 & 45.0 \ci{34.7}{55.4} & 2.17 \ci{1.47}{3.14} & 31.0 \ci{21.6}{42.0} & 0.975 & 9.2 & $-26.0$ \ci{-37.4}{-15.0} \\
MASt3R pointmap + PnP, flight height & TTA & 100 & 0.0 & 14.0 \ci{6.8}{21.6} & 10.28 \ci{8.12}{11.97} & 0.0 \ci{0.0}{0.0} & 1.04 & 43.7 & $-57.0$ \ci{-67.0}{-47.4} \\
Ours (Stage I) &  & 300 & 0.0 & 45.0 \ci{39.5}{50.7} & 1.43 \ci{1.25}{1.77} & 34.0 \ci{29.0}{39.4} & 0.720 & 4.4 & $-23.7$ \ci{-29.1}{-18.5} \\
Ours (Stage II) &  & 300 & 0.0 & 63.3 \ci{57.2}{68.9} & 0.84 \ci{0.72}{0.98} & 57.7 \ci{51.5}{63.8} & 0.666 & 1.9 & -- \\
\midrule\multicolumn{10}{@{}l}{\emph{Beyond the training range} (mean true distance 20.8\,m)}\\
SIFT + ratio test & inv. & 300 & 48.3 & 19.7 \ci{15.2}{24.0} & 18.47 \ci{17.81}{19.44} & 15.0 \ci{11.1}{18.9} & 83.6 & 90.0 & $+13.3$ \ci{+9.4}{+17.6} \\
Ours (Stage I) &  & 300 & 0.0 & 0.0 \ci{0.0}{0.0} & 18.27 \ci{17.13}{19.33} & 0.0 \ci{0.0}{0.0} & 3.88 & 48.1 & $-1.7$ \ci{-3.2}{-0.3} \\
Ours (Stage II) &  & 300 & 0.0 & 5.0 \ci{2.6}{7.4} & 18.03 \ci{16.45}{19.01} & 1.7 \ci{0.3}{3.2} & 3.32 & 45.5 & -- \\
\bottomrule
\end{tabular}}
\end{table*}


\noindent\textbf{What this experiment does not cover.} The pairs share \rlImages{} images of one site, one camera and one tilt; the intervals account for the sharing of images but not for the variation between sites. A second site, other flight heights and other illumination remain to be tested, and we do not claim more than transfer to this dataset.
